\chapter{두 번째 출력: 뇌는 어떻게 몸의 화학을 조절하는가}
\chaptermark{화학적 출력}
\label{sec:chemoutput}

\section{빠른 출력과 느린 출력}

\ref{sec:muscles}장에서 우리는 기계의 빠른 출력인 근육을 보았다. 그러나 기계에는 느린 두 번째 출력도 있다. 뇌는 체온, 심박수, 몸속 물과 당의 양을 필요한 범위 안에 유지하고, 몸을 달리기나 잠에 대비시킨다. 이를 위해 뇌는 관절을 움직이는 대신 화학 신호를 내보낸다. 경로는 두 가지다. 첫째는 내장 기관, 즉 심장, 혈관, 장, 샘으로 곧장 가는 신경이다. 둘째는 혈액이다. 일부 뉴런과 샘은 물질을 이웃 세포와의 틈이 아니라 곧바로 혈류로 내보낸다.

혈액은 이 책에 등장한 버스 가운데 가장 넓게 방송하는 버스다. \ref{sec:serocomputer}--\ref{sec:acet}장의 브로드캐스트 채널은 신호를 뇌의 넓은 영역에 보냈다. 혈액은 신호를 몸의 모든 세포에 보낸다. 혈액은 약 1분 만에 몸을 한 바퀴 돌므로, 메시지는 수십 초 안에 모든 곳에 닿고 물질이 분해될 때까지 몇 분에서 몇 시간 동안 유지된다. 이런 방송에는 주소가 전혀 없다. 명제~\ref{prop:receptor}에서처럼 메시지의 의미는 수신자가 정한다. 그 물질에 대한 수용체를 가진 세포만 반응한다.

\section{가속 페달과 브레이크: 기관으로 가는 두 가닥의 전선}

가장 좋은 예는 심장이다. 심장에는 \ref{sec:serocomputer}장의 \nt{SERO} 뉴런처럼 자체 페이스메이커가 있다. 페이스메이커는 스스로 리듬을 만들며, 신경을 모두 끊어도 심장은 1분에 약 100번 뛴다. 뇌는 리듬을 만들지 않고 조정만 한다. 그것도 부호가 반대인 두 가닥의 전선으로 한다(그림~\ref{fig:gasbrake}). 한 전선으로는 \nt{NORA}가 온다. 이것이 리듬을 빠르게 하는 \emph{가속 페달}이다. 다른 전선으로는 \nt{ACET}이 온다. 이것이 리듬을 늦추는 \emph{브레이크}다. 대략
\begin{empheq}[box=\keybox]{equation}
f \;\approx\; f_0 \;+\; a_S\,S \;-\; a_P\,P ,
\label{eq:gasbrake}
\end{empheq}
이다. 여기서 $f_0\approx100$회/분은 페이스메이커의 고유 리듬이고, $S$와 $P$는 가속 페달과 브레이크의 활동이다. 쉬고 있을 때는 브레이크가 밟혀 있어서 심장은 보통 1분에 60--70번쯤 뛴다. 운동할 때는 브레이크를 놓고 가속 페달을 밟는다.

\begin{figure}[t]
\centering
\begin{tikzpicture}[>=Stealth,
  blk/.style={draw,very thick,rounded corners=2pt,align=center,font=\small},
  pace/.style={circle,draw,very thick,double,double distance=1.2pt,minimum size=1.8cm,inner sep=0pt,font=\scriptsize,fill=red!8,align=center},
  inh/.style={-{Bar[width=7pt]},thick,red!70!black}]
\node[blk,fill=blue!5,text width=1.6cm] (B) at (0,0) {뇌의\\명령};
\node[pace] (H) at (6.2,0) {페이스\\메이커};
\draw[->,very thick,orange!85!black] (B.north east) to[bend left=20] node[above,font=\scriptsize,black,align=center] {\nt{NORA}: 가속 페달, 수 초} (H.north west);
\draw[inh,very thick,blue!60!black] (B.south east) to[bend right=20] node[below,font=\scriptsize,black,align=center] {\nt{ACET}: 브레이크, 1초 미만} (H.south west);
\draw[->,very thick] (H) -- (8.4,0) node[right,font=\small] {박동수 $f$};
\node[blk,fill=orange!10,text width=1.6cm] (A) at (0,-2.6) {\nt{ADRE}\\혈액으로};
\draw[->,thick,orange!85!black] (B) -- (A);
\foreach \y/\lab in {-2.0/{심장},-2.6/{혈관},-3.2/{간, 근육}}{
  \draw[->,thick,densely dashed,orange!85!black] (A.east) -- (4.3,\y) node[right,font=\scriptsize,black] {\lab};}
\end{tikzpicture}
\caption{심장 페이스메이커로 가는 부호가 반대인 두 가닥의 전선: 가속 페달 \nt{NORA}는 느리고 브레이크 \nt{ACET}은 빠르다. 아래는 같은 가속 페달을 브로드캐스트 버스로 보내는 경로다. \nt{ADRE} 세포가 아드레날린을 혈액으로 내보내면, 알맞은 수용체를 가진 모든 기관이 신호를 받는다(점선 화살표).}
\label{fig:gasbrake}
\end{figure}

이것은 \ref{sec:glutcomputer}장의 두 전선 부호이자 \ref{sec:muscles}장의 근육 쌍이다. 다만 이제 전선이 관절이 아니라 페이스메이커의 템포를 조절한다. 그리고 두 전선은 속도가 다르다. 둘 다 저역 통과 필터처럼 작동하지만, 브레이크는 가속 페달보다 몇 배 빠르게 변화를 통과시킨다~\cite{Berger1989}. \nt{ACET}의 경로는 짧다. 수용체와 결합한 단백질이 세포 안의 2차 전달자 없이 칼륨 채널을 직접 연다~\cite{Logothetis1987gbg}. 반면 \nt{NORA}는 세포 안의 연쇄 반응을 거쳐 작용한다. 엔지니어라면 속도가 다른 두 액추에이터를 쓰는 기법을 알아볼 것이다. 빠른 쪽은 즉각적인 미세 조정을 맡고, 느린 쪽은 크고 오래가는 변화를 맡는다.

여기서 \nt{ADRE}도 제 역할을 찾는다. 표~\ref{tab:types}에서는 뇌 안에서 \nt{ADRE}의 역할이 불분명하다고 적었다. 뇌 밖에서는 분명하다. 가속 페달 경로의 세포 일부는 기관으로 전선을 보내지 않고 아드레날린을 곧바로 혈액으로 내보낸다. 이것은 같은 가속 페달이지만 브로드캐스트 버스를 탄다. 수십 초 안에 심장, 혈관, 간, 근육에 한꺼번에 닿고 몇 분 동안 유지된다. 주소가 있는 가속 페달 \nt{NORA}는 기관 하나를 조정하고, 브로드캐스트 가속 페달 \nt{ADRE}는 몸 전체를 달리기 모드로 전환한다.

\section{되먹임이 걸린 캐스케이드}

많은 호르몬은 샘 하나가 아니라 캐스케이드가 분비한다. 뇌 바닥에 있는 작은 뉴런 무리가 첫 번째 신호를 혈액으로 내보낸다. 이 신호는 중간 샘이 두 번째 신호를 내보내게 하고, 두 번째 신호는 최종 샘이 호르몬을 내보내게 한다. 몸에 작용하는 것은 이 호르몬이다. 그리고 최종 호르몬은 앞 단계로 돌아가 그 단계들을 억제한다. 그 결과 음의 되먹임으로 감싼 3단 증폭기, 즉 식~\eqref{eq:feedback}의 \nt{SERO} 시스템과 같은 레벨 조절기가 된다.

각 단을 시상수 $\tau$와 이득 $g$를 가진 RC 회로로, 되먹임을 계수 $k$로 나타내자.
\begin{equation}
\tau\,\frac{\dd x_1}{\dd t} = -x_1 + g\bigl[u - k\,x_3\bigr]_+ ,\qquad
\tau\,\frac{\dd x_2}{\dd t} = -x_2 + g\,x_1 ,\qquad
\tau\,\frac{\dd x_3}{\dd t} = -x_3 + g\,x_2 ,
\label{eq:cascade}
\end{equation}
여기서 $u$는 뇌의 명령, $x_3$은 최종 호르몬의 농도다. 루프 이득은 $K=g^3k$이다. 평형에서 $x_3=g^3u/(1+K)$이며, 되먹임이 있는 모든 조절기처럼 루프 이득이 크면 농도가 각 단 이득의 편차에 둔감해진다. 그러나 각 단은 위상을 늦춘다. 주파수 $\omega=\sqrt3/\tau$에서 같은 단 세 개를 합치면 위상이 반 바퀴 늦어지고 진폭은 8분의 1로 줄어든다. 여기서 제어 이론의 고전적인 결과가 나온다.
\begin{empheq}[box=\keybox]{equation}
K \;<\; 8 ,
\label{eq:cascadestable}
\end{empheq}
그렇지 않으면 조절기는 주기 약 $2\pi\tau/\sqrt3\approx3.6\,\tau$로 진동하기 시작한다.

\begin{keyblock}
\begin{proposition}[정밀도 대 안정성]
\label{prop:cascade}
비례 되먹임이 걸린 3단 캐스케이드에서 농도 오차는 $1/(1+K)$이고, 캐스케이드는 $K<8$일 때만 안정하다. 따라서 이런 조절기는 농도를 약 10퍼센트보다 정확하게 유지하지 못한다. 이득을 높이려 하면 조절기는 발진기로 바뀐다.
\end{proposition}
\end{keyblock}

우리는 이를 모델~\eqref{eq:cascade}로 확인했다(그림~\ref{fig:cascade}). $K=4$에서는 켜진 뒤 농도가 조금 흔들리다가 자리를 잡는다. $K=7$에서도 자리를 잡지만 느리다. $K=12$에서는 진동이 잦아들지도 않고 끝없이 커지지도 않는다. 각 단은 음의 신호를 낼 수 없으므로, 캐스케이드는 평균 농도의 $0.83$배와 $1.24$배 사이를 오가는 일정한 맥동에 들어가며 주기는 $3.7$--$3.9\,\tau$이다.

\begin{figure}[t]
\centering
\begin{tikzpicture}[>=Stealth,x=0.3cm,y=1.2cm]
\draw[->,gray!70!black] (0,0) -- (41.5,0) node[right,font=\small,black] {$t/\tau$};
\draw[->,gray!70!black] (0,0) -- (0,2.8) node[left,font=\small,black] {$x_3/x_3^*$};
\foreach \t in {0,10,20,30,40}{\draw[gray!60!black] (\t,-0.05) -- (\t,0.05); \node[below,font=\scriptsize] at (\t,0) {\t};}
\foreach \f in {1,2}{\draw[gray!60!black] (-0.3,\f) -- (0.3,\f); \node[left,font=\scriptsize] at (0,\f) {\f};}
\draw[densely dashed,gray!60!black] (0,1) -- (40,1);
\draw[thick,blue!60!black] plot file {figures/cascade_K4.dat};
\draw[thick,red!70!black] plot file {figures/cascade_K12.dat};
\node[font=\scriptsize,blue!60!black,anchor=west] at (16,0.55) {$K=4$: 자리를 잡는다};
\node[font=\scriptsize,red!70!black,anchor=west] at (16,1.75) {$K=12$: 맥동한다};
\end{tikzpicture}
\caption{식~\eqref{eq:cascade}에 따른 되먹임이 걸린 3단 캐스케이드. 최종 호르몬의 농도를 평형값 $x_3^*$에 대한 비로 나타냈다. 루프 이득이 8보다 작으면 농도가 자리를 잡고, 8보다 크면 캐스케이드가 스스로 일정한 맥동을 만든다.}
\label{fig:cascade}
\end{figure}

실제 캐스케이드 호르몬은 정말로 맥동하며 분비된다. 예를 들어 스트레스 호르몬의 농도는 약 1시간 주기로 오르내린다. 한 가설에 따르면 이 맥동은 단계 사이의 지연된 되먹임 자체가 만들어 내며, 따로 리듬 발생기가 필요하지 않다~\cite{Walker2010}. 이는 명제~\ref{prop:ei}의 \nt{GLUT}--\nt{GABA} 루프와 조건~\eqref{eq:delaystable}의 근육 신장 루프에서 진동이 생긴 이유와 같다. 지연되는 음의 되먹임이다.

\section{호르몬의 펄스 부호}

맥동은 부작용일 뿐 아니라 부호가 되기도 한다. 생식을 조절하는 뉴런은 자신의 신호를 대략 1시간에 한 번씩 짧은 분량으로 나누어 혈액에 내보낸다. 수신 샘에 같은 신호를 분량으로 나누지 않고 연속으로 주면, 샘은 펄스를 받을 때처럼 제대로 일하지 않고 아예 멈춘다. 다시 1시간에 한 번씩 분량으로 주면 작동이 회복된다~\cite{Belchetz1978}. 즉 수신자는 농도가 아니라 분량의 \emph{빈도}를 읽는다.

엔지니어에게 여기에는 수수께끼가 없다. 물질이 오래 작용하면 지쳐서 반응을 멈추는 수용체는 \ref{sec:code}장의 고갈되는 시냅스~\eqref{eq:depression}처럼 고역 통과 필터다. 연속 신호는 막고 변화는 통과시킨다. 이런 수신자에게는 펄스로만 무언가를 전달할 수 있으므로, 정보는 농도에서 분량의 빈도와 모양으로 옮겨 간다. 게다가 분량의 빈도가 다르면 같은 샘 안의 서로 다른 세포에 다르게 작용하므로, 화학 회선 하나로 둘 이상의 값을 보낼 수 있다. \ref{sec:dofaalt}장에서 \nt{DOPA} 전선 하나로 빠른 성분과 느린 성분이 함께 전달된 것과 같다.

\section{일주기 리듬: 모드를 바꾸는 느린 발진기}

몸에서 가장 느린 리듬은 하루 주기의 리듬이다. 이 리듬은 세포 안의 지연된 음의 되먹임이 만든다. 유전자가 단백질을 만들고, 그 단백질이 몇 시간 늦게 자기 생산을 억제한다~\cite{TakahashiJS2017}. 다시 지연되는 음의 되먹임이 나온다. 이 책에서 네 번째다. 그리고 다시 리듬이 생기는데, 이번에는 주기가 약 하루다. 주 발진기는 뇌 바닥에 있는 수만 개의 뉴런으로 이루어진 작은 무리이며, 이들의 리듬이 몸의 나머지 세포를 동기화한다.

사람에서 이 발진기의 고유 주기는 정확히 하루가 아니라 평균 $24.18$시간이다~\cite{Czeisler1999}. 날마다 눈의 센서에서 오는 빛이 발진기를 맞춰 준다(\ref{sec:sensors}장). 전자 공학자에게 이것은 \emph{위상 고정 루프}다. 고유 주파수 $\omega_0$인 발진기가 주파수 $\omega$인 외부 신호에 맞춰지며, 위상차 $\varphi$는 다음 식을 따른다.
\begin{empheq}[box=\keybox]{equation}
\frac{\dd\varphi}{\dd t} \;=\; (\omega-\omega_0) \;-\; K_\varphi\,\sin\varphi .
\label{eq:pll}
\end{empheq}
$|\omega-\omega_0|<K_\varphi$이면 고정이 유지된다. 주기 $24.18$시간인 발진기는 날마다 약 11분씩 옮겨져야 하며, 보통은 아침 햇빛만으로 충분하다. 극야나 빛이 없는 동굴에서는 맞춤이 없으므로 리듬은 고유 주기로 흘러간다.

일주기 발진기는 컴퓨터의 클록 발진기가 아니다. 계산 단계를 세지도 않고 뉴런의 스파이크를 동기화하지도 않는다. 이 발진기는 \emph{모드}를 바꾼다. 잠과 깨어 있음, 호르몬 농도, 체온, 먹을 준비를 바꾼다. 이것은 \ref{sec:serocomputer}--\ref{sec:acet}장의 레지스터와 같은 종류의 느린 브로드캐스트 레지스터다. 다만 그 값이 태양에 맞춘 일정표에 따라 바뀐다.

\begin{keyblock}
\begin{proposition}[화학적 출력]
\label{prop:chemoutput}
기계의 느린 출력은 빠른 출력과 같은 부품으로 만들어진다. 기관은 부호가 반대이고 속도가 다른 두 가닥의 전선, 즉 가속 페달 \nt{NORA}와 브레이크 \nt{ACET}을 받고, 몸 전체는 아드레날린 \nt{ADRE}를 비롯한 혈액의 브로드캐스트 신호를 한꺼번에 받는다. 호르몬 농도는 음의 되먹임이 걸린 캐스케이드가 유지하며, 그 정밀도는 안정성에 제한된다. 일부 호르몬은 분량의 빈도로 부호화된다. 일주기 리듬은 빛에 위상 고정되는 느린 발진기다. 경로의 구조, 분량 실험, 주기 실험은 측정되었다. 캐스케이드의 맥동을 되먹임 자체로 설명하는 것은 가설이다.
\end{proposition}
\end{keyblock}

\section{요약}

\begin{table}[t]
\centering
\footnotesize
\setlength{\tabcolsep}{4pt}
\renewcommand{\arraystretch}{1.15}
\begin{tabular}{>{\raggedright}p{3.1cm}>{\raggedright}p{4.8cm}>{\raggedright\arraybackslash}p{4.9cm}}
\toprule
화학적 출력이 하는 일 & 방법 & 알려진 것 \\
\midrule
기관을 조정한다 & 가속 페달 \nt{NORA}와 브레이크 \nt{ACET}~\eqref{eq:gasbrake} & 측정됨; 브레이크가 가속 페달보다 빠르다 \\
몸 전체를 달리기 모드로 바꾼다 & 혈액 속 아드레날린 \nt{ADRE} & 측정됨 \\
호르몬 농도를 유지한다 & 되먹임이 걸린 캐스케이드, $K<8$~\eqref{eq:cascadestable} & 캐스케이드와 되먹임은 측정됨; 루프 자체가 만드는 맥동은 가설 \\
분량의 빈도로 전달한다 & 고역 통과 필터로서의 지친 수용체 & 분량 의존성은 측정됨 \\
하루 동안 모드를 바꾼다 & 빛에 위상 고정되는 발진기~\eqref{eq:pll} & 주기와 빛에 의한 맞춤은 측정됨 \\
\bottomrule
\end{tabular}
\caption{기계는 몸의 화학을 어떻게 조절하는가.}
\label{tab:chemoutput}
\end{table}

기계에는 출력이 두 개 있다(표~\ref{tab:chemoutput}). 빠른 출력은 근육이다. 밀리초 단위이고, 주소가 있으며, 관절 하나하나로 간다. 느린 출력은 몸의 화학이다. 초, 분, 하루 단위이고, 두 가닥의 전선으로 기관에, 혈액으로 모든 곳에 한꺼번에 간다. 두 출력의 부품은 같다. 부호가 반대인 두 전선, 페이스메이커, 되먹임과 그 지연이다. 기계 내부를 이루는 부품도 바로 이것들이다.

\section{연습문제}

\begin{exercise}[\lvlA]
\label{ex:16:1}
\emph{필터로서의 혈액.} 호르몬은 반감기 $t_{1/2}=10$~min으로 혈액에서 제거되고, $T=60$~min마다 짧은 분량으로 분비된다. 농도의 최댓값은 최솟값의 몇 배이고, 분량의 기본 고조파는 몇 배로 감쇠하는가? $t_{1/2}$가 얼마일 때 최댓값이 최솟값의 2배에 불과한가?
\end{exercise}

\begin{exercise}[\lvlA]
\label{ex:16:2}
\emph{일주기 위상 고정.} 발진기~\eqref{eq:pll}의 고유 주기는 $24.18$~h이고, 빛은 위상을 하루에 최대 1시간 옮긴다: $K_\varphi=2\pi/24$~rad/day. 정상 상태 위상차 $\varphi^*$를 분 단위로, 그리고 이완 시간을 구하라. 외부 신호가 $\pm6$~h 도약한 뒤 불일치가 1시간 아래로 줄어드는 데 며칠이 걸리는가?
\end{exercise}

\begin{exercise}[\lvlB]
\label{ex:16:3}
\emph{정밀도 대 안정성.} 선형화한 캐스케이드~\eqref{eq:cascade}를 같은 단 $n$개로 늘렸다. 임계 이득 $K_c(n)$과, $n=3$과 $n=6$에서 달성할 수 있는 최소 오차 $1/(1+K_c)$를 구하라. $n\to\infty$일 때 이 오차는 어디로 수렴하는가?
\end{exercise}

\begin{exercise}[\lvlB]
\label{ex:16:4}
\emph{가속 페달과 브레이크.} 식~\eqref{eq:gasbrake}에서 가속 페달과 브레이크의 기여는 $\tau_S=2$~s와 $\tau_P=0.4$~s인 RC 회로의 출력이고, 명령은 각각 분당 $80$회와 $60$회를 넘지 않는다. $f_0=100$이고, 안정 시 브레이크 $35$, 가속 페달 $0$, $f=65$이다. $f=120$에 도달하는 데 얼마나 걸리는가? 브레이크를 놓지 않고 가속 페달만 쓰면 어떤가?
\end{exercise}

\begin{exercise}[\lvlB]
\label{ex:16:5}
\emph{모델링: 지연이 있는 캐스케이드.} \texttt{models/\allowbreak chem\_models.py}의 함수 \texttt{cascade}를 복사해(파일은 실행하지 말 것) 되먹임에 지연 $d$를 넣어라. 첫 단이 $x_3(t-d)$를 본다. $d/\tau=0$, $0.5$, $1$, $2$에서 $K_c$와 경계에서의 주기를 구하고, $0.9K_c$와 $1.1K_c$에서 시뮬레이션으로 확인하라.
\end{exercise}

\begin{exercise}[\lvlB]
\label{ex:16:6}
\emph{분량을 읽는 수신자.} 수용체가 지친다: $y=[c-a]_+$, $\tau_a\,\dd a/\dd t=c-a$, $\tau_a=30$~min. 호르몬은 평균 용량 $\bar c$는 같게, 한 번에 $6$~min씩 주기 $T$로 들어온다. $T=15$, $60$, $120$~min과 $T\to\infty$에서 평균 반응을 구하라. 농도가 $\bar c$로 일정하면 반응은 얼마인가?
\end{exercise}

\begin{exercise}[\lvlC]
\label{ex:16:7}
\emph{발진기인가, 되먹임인가?} 캐스케이드~\eqref{eq:cascade}의 지연된 되먹임이 만드는 맥동과 별도의 페이스메이커가 만드는 맥동을 구별하는 실험을 제안하라. $K$를 1.5배밖에 바꿀 수 없고 주기를 $5\%$ 정확도로 측정한다면, 주기의 이동으로 두 가설을 구별할 수 있는가?
\end{exercise}
